\chapter{Advanced Traits and Generics}
\label{ch:traits-generics}

\begin{quote}
\textit{``Abstraction is not about vagueness, it is about being precise on a new semantic level.''} --- Edsger W. Dijkstra
\end{quote}

Chapter 4 introduced traits as a way to define shared behavior across types. This chapter goes deeper, covering operator overloading, the complete trait ecosystem, advanced generic patterns, and real-world idioms that make Novus code both flexible and efficient.

By the end of this chapter, you'll understand how to:
\begin{itemize}
    \item Make your types work with Novus operators (\texttt{+}, \texttt{-}, \texttt{*}, \texttt{[]}, etc.)
    \item Implement comparison and ordering for custom types
    \item Create fallible and infallible type conversions
    \item Build custom iterators for any data structure
    \item Write generic numeric algorithms
    \item Design flexible APIs with reference traits
    \item Apply common patterns like newtype, extension traits, and blanket implementations
\end{itemize}

\section{Operator Overloading}
\label{sec:operator-overloading}

Novus allows you to define how operators work with your custom types by implementing operator traits. When you implement the \texttt{Add} trait for a type, the \texttt{+} operator works with values of that type.

\subsection{Arithmetic Operator Traits}

The arithmetic operator traits are defined in \texttt{std::core}:

\begin{lstlisting}
pub trait Add {
    fn add(&self, other: &Self) -> Self
}

pub trait Sub {
    fn sub(&self, other: &Self) -> Self
}

pub trait Mul {
    fn mul(&self, other: &Self) -> Self
}

pub trait Div {
    fn div(&self, other: &Self) -> Self
}

pub trait Rem {
    fn rem(&self, other: &Self) -> Self
}

pub trait Neg {
    fn neg(&self) -> Self
}
\end{lstlisting}

Each trait method takes references and returns an owned value. This design avoids unnecessary copies while keeping the interface simple.

\subsubsection{Complete Operator-to-Trait Mapping}

\begin{table}[h]
\centering
\begin{tabular}{lll}
\hline
\textbf{Operator} & \textbf{Trait} & \textbf{Method Signature} \\
\hline
\texttt{a + b} & \texttt{Add} & \texttt{fn add(\&self, other: \&Self) -> Self} \\
\texttt{a - b} & \texttt{Sub} & \texttt{fn sub(\&self, other: \&Self) -> Self} \\
\texttt{a * b} & \texttt{Mul} & \texttt{fn mul(\&self, other: \&Self) -> Self} \\
\texttt{a / b} & \texttt{Div} & \texttt{fn div(\&self, other: \&Self) -> Self} \\
\texttt{a \% b} & \texttt{Rem} & \texttt{fn rem(\&self, other: \&Self) -> Self} \\
\texttt{-a} & \texttt{Neg} & \texttt{fn neg(\&self) -> Self} \\
\hline
\end{tabular}
\caption{Arithmetic Operator Traits}
\end{table}

\subsection{Implementing Arithmetic Operators}

Let's implement arithmetic for a 2D vector type:

\begin{lstlisting}
from std::core import Add, Sub, Neg

struct Vec2 {
    x: i32,
    y: i32,
}

impl Vec2 {
    pub fn new(x: i32, y: i32) -> Vec2 {
        return Vec2 { x: x, y: y }
    }
}

impl Add for Vec2 {
    fn add(&self, other: &Vec2) -> Vec2 {
        return Vec2 {
            x: self.x + other.x,
            y: self.y + other.y
        }
    }
}

impl Sub for Vec2 {
    fn sub(&self, other: &Vec2) -> Vec2 {
        return Vec2 {
            x: self.x - other.x,
            y: self.y - other.y
        }
    }
}

impl Neg for Vec2 {
    fn neg(&self) -> Vec2 {
        return Vec2 { x: -self.x, y: -self.y }
    }
}
\end{lstlisting}

Now you can use operators naturally:

\begin{lstlisting}
// ... (with Vec2 and operator impls defined above)
let a = Vec2::new(10, 20)
let b = Vec2::new(5, 15)

let sum = a + b       // Vec2 { x: 15, y: 35 }
let diff = a - b      // Vec2 { x: 5, y: 5 }
let negated = -a      // Vec2 { x: -10, y: -20 }
\end{lstlisting}

\subsection{Fixed-Point Arithmetic Example}

Here's how \texttt{fixed32} (16.16 fixed-point) implements multiplication:

\begin{lstlisting}
// ... (illustrative - actual impl is in compiler)
impl Mul for fixed32 {
    fn mul(&self, other: &fixed32) -> fixed32 {
        // Fixed-point multiplication requires shifting
        // a * b = (a_raw * b_raw) >> 16
        // This is handled by the compiler intrinsic
        return fixed32_mul(*self, *other)
    }
}
\end{lstlisting}

\textbf{Performance note:} On 68000, 32-bit multiplication is done in software. The 68020+ have hardware \texttt{MULU.L} and \texttt{MULS.L} instructions. The compiler generates appropriate code based on your \texttt{--cpu} target.

\subsection{Common Mistakes}

\textbf{Mistake 1: Forgetting that methods take references}

\begin{lstlisting}
// WRONG - takes self by value
impl Add for Vec2 {
    fn add(self, other: Vec2) -> Vec2 {  // Wrong signature!
        // ...
    }
}

// CORRECT - takes references
impl Add for Vec2 {
    fn add(&self, other: &Vec2) -> Vec2 {
        // ...
    }
}
\end{lstlisting}

\textbf{Mistake 2: Division by zero}

Operator traits don't have built-in error handling. If your type can encounter division by zero, you have options:

\begin{lstlisting}
// Option 1: Panic (simple but crashes)
impl Div for Fraction {
    fn div(&self, other: &Fraction) -> Fraction {
        if other.numerator == 0 {
            panic!("Division by zero")
        }
        // ... normal division
    }
}

// Option 2: Saturate or return a sentinel value
impl Div for Fixed32 {
    fn div(&self, other: &Fixed32) -> Fixed32 {
        if other.is_zero() {
            return Fixed32::max_value()  // Saturate to max
        }
        // ... normal division
    }
}
\end{lstlisting}

For fallible operations, consider providing a separate \texttt{try\_div} method that returns \texttt{Result}.

\subsection{Index Operators}
\label{sec:index-operators}

The \texttt{Index} and \texttt{IndexMut} traits enable the \texttt{[]} operator for your types.

\begin{lstlisting}
pub trait Index<I, T> {
    fn index(&self, idx: I) -> T
}

pub trait IndexMut<I, T> {
    fn index_set(&var self, idx: I, value: T)
}
\end{lstlisting}

Note the asymmetry: \texttt{Index} returns a value, while \texttt{IndexMut} takes a value to set. This differs from some other languages where both return references.

\subsubsection{Implementing Index for a Grid}

\begin{lstlisting}
from std::core import Index, IndexMut

struct Grid {
    data: Vec<i32>,
    width: u32,
    height: u32,
}

impl Grid {
    pub fn new(width: u32, height: u32) -> Option<Grid> {
        let size = width * height
        match Vec::with_capacity(size) {
            Option::Some(var v) => {
                // Initialize with zeros
                for i in 0..size {
                    v.push(0)
                }
                return Option::Some(Grid {
                    data: v,
                    width: width,
                    height: height
                })
            },
            Option::None => return Option::None
        }
    }

    fn calc_index(&self, x: u32, y: u32) -> u32 {
        return y * self.width + x
    }
}

// Read access: let value = grid[(x, y)]
impl Index<(u32, u32), i32> for Grid {
    fn index(&self, pos: (u32, u32)) -> i32 {
        let (x, y) = pos
        let idx = self.calc_index(x, y)
        return self.data.get(idx).unwrap()
    }
}

// Write access: grid[(x, y)] = value
impl IndexMut<(u32, u32), i32> for Grid {
    fn index_set(&var self, pos: (u32, u32), value: i32) {
        let (x, y) = pos
        let idx = self.calc_index(x, y)
        self.data.set(idx, value)
    }
}
\end{lstlisting}

Usage:

\begin{lstlisting}
// ... (with Grid defined above)
var grid = Grid::new(10, 10).unwrap()

grid[(5, 3)] = 42           // Write using IndexMut
let value = grid[(5, 3)]    // Read using Index
\end{lstlisting}

\subsubsection{Bounds Checking}

In debug builds, you should add bounds checking to catch programming errors:

\begin{lstlisting}
// ... (with Grid defined above)
impl Index<(u32, u32), i32> for Grid {
    fn index(&self, pos: (u32, u32)) -> i32 {
        let (x, y) = pos
        if x >= self.width || y >= self.height {
            panic!("Grid index out of bounds")
        }
        let idx = self.calc_index(x, y)
        return self.data.get(idx).unwrap()
    }
}
\end{lstlisting}

For production code where bounds are guaranteed, you can use unchecked access for performance.


\section{Comparison and Equality Traits}
\label{sec:comparison-traits}

Beyond the basic \texttt{Eq} trait (covered in Chapter 4), Novus provides \texttt{PartialOrd} and \texttt{Ord} for ordering comparisons.

\subsection{The Equality Contract}

The \texttt{Eq} trait requires these mathematical properties:

\begin{itemize}
    \item \textbf{Reflexive}: \texttt{a.eq(\&a)} must return \texttt{true}
    \item \textbf{Symmetric}: \texttt{a.eq(\&b)} must equal \texttt{b.eq(\&a)}
    \item \textbf{Transitive}: if \texttt{a.eq(\&b)} and \texttt{b.eq(\&c)}, then \texttt{a.eq(\&c)}
\end{itemize}

Violating these properties will cause collections like \texttt{HashMap} to behave incorrectly.

\subsection{PartialOrd --- Comparison Operators}

\texttt{PartialOrd} enables the \texttt{<}, \texttt{<=}, \texttt{>}, and \texttt{>=} operators:

\begin{lstlisting}
pub trait PartialOrd {
    fn lt(&self, other: &Self) -> bool   // <
    fn le(&self, other: &Self) -> bool   // <=
    fn gt(&self, other: &Self) -> bool   // >
    fn ge(&self, other: &Self) -> bool   // >=
}
\end{lstlisting}

\subsubsection{Implementing PartialOrd}

\begin{lstlisting}
from std::core import Eq, PartialOrd

struct Version {
    major: u32,
    minor: u32,
    patch: u32,
}

impl Eq for Version {
    fn eq(&self, other: &Version) -> bool {
        if self.major != other.major { return false }
        if self.minor != other.minor { return false }
        return self.patch == other.patch
    }
}

impl PartialOrd for Version {
    fn lt(&self, other: &Version) -> bool {
        if self.major != other.major {
            return self.major < other.major
        }
        if self.minor != other.minor {
            return self.minor < other.minor
        }
        return self.patch < other.patch
    }

    fn le(&self, other: &Version) -> bool {
        return self.lt(other) || self.eq(other)
    }

    fn gt(&self, other: &Version) -> bool {
        return other.lt(self)
    }

    fn ge(&self, other: &Version) -> bool {
        return other.le(self)
    }
}
\end{lstlisting}

Usage:

\begin{lstlisting}
// ... (with Version defined above)
let v1 = Version { major: 1, minor: 2, patch: 0 }
let v2 = Version { major: 1, minor: 3, patch: 0 }

if v1 < v2 {
    println("v1 is older than v2")
}

if v1 >= v2 {
    println("v1 is at least as new as v2")
}
\end{lstlisting}

\subsection{Ord --- Total Ordering}

\texttt{Ord} is for types where every value is comparable. It's required for sorting and binary search.

\begin{lstlisting}
pub trait Ord {
    /// Returns: negative if self < other, zero if equal, positive if self > other
    fn cmp(&self, other: &Self) -> i32
}
\end{lstlisting}

The return value convention:
\begin{itemize}
    \item Negative (\texttt{< 0}): \texttt{self} is less than \texttt{other}
    \item Zero (\texttt{== 0}): \texttt{self} equals \texttt{other}
    \item Positive (\texttt{> 0}): \texttt{self} is greater than \texttt{other}
\end{itemize}

\begin{lstlisting}
from std::core import Ord

// ... (with Version and Eq defined above)
impl Ord for Version {
    fn cmp(&self, other: &Version) -> i32 {
        if self.major != other.major {
            return (i32)self.major - (i32)other.major
        }
        if self.minor != other.minor {
            return (i32)self.minor - (i32)other.minor
        }
        return (i32)self.patch - (i32)other.patch
    }
}
\end{lstlisting}

With \texttt{Ord} implemented, your type works with sorting:

\begin{lstlisting}
// ... (with Version and Ord defined above)
var versions: Vec<Version> = load_versions()
versions.sort()  // Sorts using Ord::cmp
\end{lstlisting}

\subsection{When to Use PartialOrd vs Ord}

\begin{itemize}
    \item Use \texttt{PartialOrd} when comparison might not always be possible (e.g., floating-point with NaN)
    \item Use \texttt{Ord} when all values are comparable (integers, strings, versions)
\end{itemize}

Floating-point types implement \texttt{PartialOrd} but not \texttt{Ord} because NaN is not comparable to anything, including itself.

\subsection{Deriving Comparison Traits}

For simple structs, use \texttt{@derive}:

\begin{lstlisting}
#[derive(Eq)]
struct Point {
    x: i32,
    y: i32,
}
\end{lstlisting}

The derived \texttt{Eq} compares all fields in declaration order. Note that \texttt{@derive} currently supports \texttt{Eq} and \texttt{Hash}, but not \texttt{Ord}.

\subsection{Common Mistakes}

\textbf{Mistake 1: Inconsistent Eq and Hash}

If two values are equal via \texttt{Eq}, they \emph{must} produce the same hash:

\begin{lstlisting}
// WRONG - case-insensitive eq but case-sensitive hash
impl Eq for CaseInsensitiveString {
    fn eq(&self, other: &Self) -> bool {
        return self.to_lowercase() == other.to_lowercase()
    }
}

impl Hash for CaseInsensitiveString {
    fn hash(&self) -> u32 {
        return self.data.hash()  // BUG: "Hello" and "hello" hash differently!
    }
}

// CORRECT - both use lowercase
impl Hash for CaseInsensitiveString {
    fn hash(&self) -> u32 {
        return self.to_lowercase().hash()
    }
}
\end{lstlisting}

\textbf{Mistake 2: Ord that disagrees with Eq}

If \texttt{a.cmp(\&b) == 0}, then \texttt{a.eq(\&b)} should return \texttt{true}. Violating this causes subtle bugs in sorted collections.


\section{Conversion Traits}
\label{sec:conversion-traits}

Novus provides a family of traits for type conversion, covered in Chapter 4's \texttt{From} section. Here we cover the complete picture.

\subsection{The Conversion Trait Family}

\begin{table}[h]
\centering
\begin{tabular}{llp{6cm}}
\hline
\textbf{Trait} & \textbf{Method} & \textbf{Use Case} \\
\hline
\texttt{From<T>} & \texttt{convert(T) -> Self} & Infallible conversion from T \\
\texttt{Into<T>} & \texttt{into(self) -> T} & Infallible conversion to T \\
\texttt{TryFrom<T, E>} & \texttt{try\_from(T) -> Result<Self, E>} & Fallible conversion from T \\
\texttt{TryInto<T, E>} & \texttt{try\_into(self) -> Result<T, E>} & Fallible conversion to T \\
\hline
\end{tabular}
\caption{Conversion Traits}
\end{table}

\subsection{From and Into}

\texttt{From<T>} was covered in Chapter 4. The key insight is that \texttt{Into} is the reverse direction:

\begin{lstlisting}
pub trait From<T> {
    fn convert(value: T) -> Self
}

pub trait Into<T> {
    fn into(consuming self) -> T
}
\end{lstlisting}

If you implement \texttt{From<A>} for B, Novus can automatically provide \texttt{Into<B>} for A. Always prefer implementing \texttt{From} over \texttt{Into}.

\subsection{TryFrom and TryInto}

When conversion can fail, use the fallible variants:

\begin{lstlisting}
pub trait TryFrom<T, E> {
    fn try_from(value: T) -> Result<Self, E>
}

pub trait TryInto<T, E> {
    fn try_into(consuming self) -> Result<T, E>
}
\end{lstlisting}

\subsubsection{Example: Parsing with TryFrom}

\begin{lstlisting}
from std::core import TryFrom, Result

pub enum ParseError {
    InvalidFormat,
    OutOfRange,
    Empty,
}

impl TryFrom<*u8, ParseError> for u8 {
    fn try_from(s: *u8) -> Result<u8, ParseError> {
        if s == null {
            return Result::Err(ParseError::Empty)
        }

        // Simple ASCII digit parsing
        var value: u32 = 0
        var ptr = s
        while *ptr != 0 {
            let c = *ptr
            if c < '0' || c > '9' {
                return Result::Err(ParseError::InvalidFormat)
            }
            value = value * 10 + (u32)(c - '0')
            if value > 255 {
                return Result::Err(ParseError::OutOfRange)
            }
            ptr = ptr + 1
        }
        return Result::Ok((u8)value)
    }
}
\end{lstlisting}

Usage:

\begin{lstlisting}
// ... (with u8 TryFrom impl defined above)
match u8::try_from("255") {
    Result::Ok(n) => println(f"Parsed: {n}"),
    Result::Err(ParseError::OutOfRange) => println("Number too large"),
    Result::Err(ParseError::InvalidFormat) => println("Not a number"),
    Result::Err(ParseError::Empty) => println("Empty string"),
}
\end{lstlisting}

\subsection{Designing Error Types for Conversions}

Good conversion error types are:
\begin{itemize}
    \item \textbf{Specific}: Different failure modes get different variants
    \item \textbf{Informative}: Include context when useful
    \item \textbf{Lightweight}: Don't allocate in the error path if possible
\end{itemize}

\begin{lstlisting}
// Good - specific variants, no allocation
pub enum NumericParseError {
    Empty,
    InvalidCharacter,
    Overflow,
    Underflow,
}

// Less good - loses information
pub enum GenericError {
    Failed,
}
\end{lstlisting}

\subsection{Choosing the Right Conversion Trait}

\begin{itemize}
    \item Use \texttt{From}/\texttt{Into} when conversion always succeeds
    \item Use \texttt{TryFrom}/\texttt{TryInto} when conversion can fail
    \item Implement \texttt{From}, not \texttt{Into} (you get \texttt{Into} for free)
    \item Never panic in \texttt{From} --- if it can fail, use \texttt{TryFrom}
\end{itemize}


\section{The Iterator System}
\label{sec:iterators}

Iterators provide a uniform way to traverse sequences. Implementing \texttt{Iterator} lets your type work with loops and iterator-consuming functions.

\subsection{The Iterator Trait}

\begin{lstlisting}
pub trait Iterator<T> {
    fn next(&var self) -> Option<T>
}
\end{lstlisting}

The contract:
\begin{itemize}
    \item \texttt{next()} returns \texttt{Some(value)} for each element
    \item When exhausted, returns \texttt{None}
    \item After returning \texttt{None}, should continue returning \texttt{None}
\end{itemize}

\subsection{Creating a Custom Iterator}

Let's implement an iterator for a simple range:

\begin{lstlisting}
from std::core import Iterator, Option

struct RangeIter {
    current: i32,
    end: i32,
}

impl RangeIter {
    pub fn new(start: i32, end: i32) -> RangeIter {
        return RangeIter { current: start, end: end }
    }
}

impl Iterator<i32> for RangeIter {
    fn next(&var self) -> Option<i32> {
        if self.current >= self.end {
            return Option::None
        }
        let value = self.current
        self.current = self.current + 1
        return Option::Some(value)
    }
}
\end{lstlisting}

Usage:

\begin{lstlisting}
// ... (with RangeIter defined above)
var iter = RangeIter::new(0, 5)
loop {
    match iter.next() {
        Option::Some(n) => println(f"{n}"),
        Option::None => break,
    }
}
// Output: 0, 1, 2, 3, 4
\end{lstlisting}

\subsection{Iterating Collections}

The \texttt{Iterable} trait connects collections to for-loops:

\begin{lstlisting}
pub trait Iterable<T> {
    fn get(&self, index: u32) -> Option<T>
    fn len(&self) -> u32
}
\end{lstlisting}

When you write \texttt{for item in collection}, the compiler uses \texttt{len()} to determine bounds and \texttt{get()} to access each element by index.

\subsubsection{Vec Iterator Example}

\begin{lstlisting}
// ... (illustrative - actual impl in stdlib)
struct VecIter<T> {
    vec: *Vec<T>,
    index: u32,
}

impl<T> Iterator<T> for VecIter<T> {
    fn next(&var self) -> Option<T> {
        if self.index >= (*self.vec).len() {
            return Option::None
        }
        let value = (*self.vec).get(self.index)
        self.index = self.index + 1
        return value
    }
}
\end{lstlisting}

\subsection{Iterator for a Linked List}

Here's a more complex example --- iterating a singly-linked list:

\begin{lstlisting}
from std::core import Iterator, Option

struct Node<T> {
    value: T,
    next: *Node<T>,
}

struct LinkedList<T> {
    head: *Node<T>,
    len: u32,
}

struct LinkedListIter<T> {
    current: *Node<T>,
}

impl<T> LinkedList<T> {
    pub fn iter(&self) -> LinkedListIter<T> {
        return LinkedListIter { current: self.head }
    }
}

impl<T> Iterator<T> for LinkedListIter<T> {
    fn next(&var self) -> Option<T> {
        if self.current == null {
            return Option::None
        }
        let value = (*self.current).value
        self.current = (*self.current).next
        return Option::Some(value)
    }
}
\end{lstlisting}

\subsection{Performance Considerations}

Iterators in Novus are designed for zero overhead:

\begin{itemize}
    \item Iterator structs are typically small (a pointer and an index)
    \item \texttt{next()} calls are inlined by the optimizer
    \item No heap allocation for iteration itself
    \item Performance matches hand-written loops
\end{itemize}

When iterators might be slower:
\begin{itemize}
    \item Very tight loops where function call overhead matters
    \item When the optimizer can't inline (rare)
    \item Complex iterator chains (consider breaking into loops)
\end{itemize}

\subsection{Common Mistakes}

\textbf{Mistake 1: Modifying the collection while iterating}

\begin{lstlisting}
// ... (illustrative - shows incorrect pattern)
// WRONG - undefined behavior
var vec = Vec::new()
vec.push(1)
vec.push(2)
vec.push(3)

var iter = vec.iter()
loop {
    match iter.next() {
        Option::Some(n) => {
            if n == 2 {
                vec.push(4)  // BUG: modifying while iterating!
            }
        },
        Option::None => break,
    }
}
\end{lstlisting}

\textbf{Mistake 2: Using an exhausted iterator}

\begin{lstlisting}
// ... (illustrative - shows incorrect pattern)
var iter = vec.iter()
// Consume all elements
loop {
    match iter.next() {
        Option::Some(_) => {},
        Option::None => break,
    }
}

// This won't iterate anything - iterator is exhausted
loop {
    match iter.next() {
        Option::Some(n) => println(f"{n}"),  // Never reached
        Option::None => break,
    }
}
\end{lstlisting}

Create a new iterator if you need to iterate again.


\section{Numeric Traits}
\label{sec:numeric-traits}

These traits enable generic programming over numeric types --- essential for math libraries and algorithms.

\subsection{Zero and One}

\begin{lstlisting}
pub trait Zero {
    fn zero() -> Self
    fn is_zero(&self) -> bool
}

pub trait One {
    fn one() -> Self
}
\end{lstlisting}

These represent the additive and multiplicative identity elements.

\subsubsection{Implementations for Primitive Types}

All numeric primitives implement these traits:

\begin{lstlisting}
// ... (implemented in std::core)
impl Zero for i32 {
    fn zero() -> i32 { return 0 }
    fn is_zero(&self) -> bool { return *self == 0 }
}

impl One for i32 {
    fn one() -> i32 { return 1 }
}
\end{lstlisting}

\subsection{Writing Generic Numeric Functions}

\begin{lstlisting}
from std::core import Zero, One, Add, Mul

fn sum<T>(items: *T, count: u32) -> T where T: Zero + Add {
    var total = T::zero()
    for i in 0..count {
        total = total + *(items + i)
    }
    return total
}

fn product<T>(items: *T, count: u32) -> T where T: One + Mul {
    var result = T::one()
    for i in 0..count {
        result = result * *(items + i)
    }
    return result
}
\end{lstlisting}

Usage:

\begin{lstlisting}
// ... (with sum and product defined above)
let numbers = [1, 2, 3, 4, 5]
let total = sum(&numbers[0], 5)      // 15
let prod = product(&numbers[0], 5)   // 120
\end{lstlisting}

\subsection{Bounded --- Min and Max Values}

\begin{lstlisting}
pub trait Bounded {
    fn min_value() -> Self
    fn max_value() -> Self
}
\end{lstlisting}

\subsubsection{Bounds for Novus Numeric Types}

\begin{table}[h]
\centering
\begin{tabular}{lrr}
\hline
\textbf{Type} & \textbf{min\_value()} & \textbf{max\_value()} \\
\hline
\texttt{i8} & $-128$ & $127$ \\
\texttt{i16} & $-32{,}768$ & $32{,}767$ \\
\texttt{i32} & $-2{,}147{,}483{,}648$ & $2{,}147{,}483{,}647$ \\
\texttt{u8} & $0$ & $255$ \\
\texttt{u16} & $0$ & $65{,}535$ \\
\texttt{u32} & $0$ & $4{,}294{,}967{,}295$ \\
\hline
\end{tabular}
\caption{Numeric Type Bounds}
\end{table}

\subsection{Practical Use: Saturating Arithmetic}

Saturating arithmetic clamps at bounds instead of overflowing:

\begin{lstlisting}
// ... (illustrative - generic numeric functions)
from std::core import Add, Bounded, PartialOrd

fn saturating_add<T>(a: T, b: T) -> T where T: Add + Bounded + PartialOrd {
    let result = a + b
    // Check for overflow: if result wrapped around
    if result.lt(&a) {
        return T::max_value()
    }
    return result
}

fn clamp<T>(value: T, min: T, max: T) -> T where T: PartialOrd {
    if value.lt(&min) { return min }
    if value.gt(&max) { return max }
    return value
}
\end{lstlisting}

\subsection{Fixed-Point and Bounded}

Fixed-point types also implement \texttt{Bounded}:

\begin{lstlisting}
// ... (illustrative - actual values depend on representation)
impl Bounded for fixed32 {
    fn min_value() -> fixed32 {
        // -32768.0 in 16.16 format
        return fixed32_from_raw(-2147483648)
    }
    fn max_value() -> fixed32 {
        // 32767.99998... in 16.16 format
        return fixed32_from_raw(2147483647)
    }
}
\end{lstlisting}


\section{Reference Traits}
\label{sec:reference-traits}

\texttt{AsRef} and \texttt{AsMut} enable functions to accept multiple types that can be ``viewed as'' a common type.

\subsection{AsRef --- Cheap Reference Conversion}

\begin{lstlisting}
pub trait AsRef<T> {
    fn as_ref(&self) -> T
}
\end{lstlisting}

Note: In Novus, \texttt{AsRef} returns \texttt{T} directly (typically a slice or reference type), not \texttt{\&T}.

\subsection{Designing Flexible APIs}

Without \texttt{AsRef}:

\begin{lstlisting}
// Inflexible - only accepts Str
fn count_lines(s: Str) -> u32 {
    var count: u32 = 0
    for c in s.bytes() {
        if c == '\n' {
            count = count + 1
        }
    }
    return count
}
\end{lstlisting}

With \texttt{AsRef}:

\begin{lstlisting}
// Flexible - accepts String, Str, or anything AsRef<Str>
fn count_lines<S>(s: &S) -> u32 where S: AsRef<Str> {
    let str_view = s.as_ref()
    var count: u32 = 0
    // ... count newlines
    return count
}
\end{lstlisting}

\subsection{AsMut --- Mutable View Conversion}

\begin{lstlisting}
pub trait AsMut<T> {
    fn as_mut(&var self) -> T
}
\end{lstlisting}

Use when your function needs to modify the underlying data.

\subsection{Standard Library Usage}

Many stdlib functions use \texttt{AsRef} for flexibility:

\begin{lstlisting}
// ... (illustrative - file opening API)
pub fn open<P>(path: &P, mode: u32) -> Result<FileHandle, DosError>
    where P: AsRef<Str>
{
    let path_str = path.as_ref()
    // ... open the file using path_str
}
\end{lstlisting}

Callers benefit:

\begin{lstlisting}
// Works with string literal (Str)
let f1 = open(&"config.txt", MODE_OLDFILE)

// Works with owned String
var path = String::new()
path.push_str("data/file.dat")
let f2 = open(&path, MODE_OLDFILE)
\end{lstlisting}

\subsection{When Not to Use AsRef}

\begin{itemize}
    \item When you need ownership (use \texttt{Into} instead)
    \item When the conversion is expensive (AsRef should be cheap)
    \item When it makes the API confusing (simpler is sometimes better)
\end{itemize}


\section{Advanced Generics}
\label{sec:advanced-generics}

This section covers more sophisticated generic patterns.

\subsection{Multiple Trait Bounds}

Use \texttt{+} to require multiple traits:

\begin{lstlisting}
fn debug_and_hash<T>(value: &T) -> u32 where T: Debug + Hash {
    println(value.debug_string())
    return value.hash()
}
\end{lstlisting}

\subsection{Multiple Type Parameters}

\begin{lstlisting}
// ... (illustrative - multiple generic parameters)
fn convert_and_store<T, U>(input: T, storage: &var Vec<U>) where T: Into<U> {
    let converted: U = input.into()
    storage.push(converted)
}
\end{lstlisting}

\subsection{Where Clauses on Structs}

Constrain which types can instantiate a generic struct:

\begin{lstlisting}
// ... (illustrative - where clause on struct)
struct SortedVec<T> where T: Ord {
    items: Vec<T>,
}

impl<T> SortedVec<T> where T: Ord {
    pub fn new() -> Option<SortedVec<T>> {
        match Vec::new() {
            Option::Some(v) => {
                return Option::Some(SortedVec { items: v })
            },
            Option::None => return Option::None
        }
    }

    pub fn insert(&var self, value: T) {
        // Find insertion point to maintain sorted order
        var idx: u32 = 0
        for i in 0..self.items.len() {
            match self.items.get(i) {
                Option::Some(item) => {
                    if item.cmp(&value) > 0 {
                        break
                    }
                },
                Option::None => break,
            }
            idx = i + 1
        }
        self.items.insert(idx, value)
    }
}
\end{lstlisting}

\subsection{Generic Implementations}

Implement traits generically for families of types:

\begin{lstlisting}
// ... (illustrative - Display for all Option<T> where T: Display)
impl<T> Display for Option<T> where T: Display {
    fn fmt(&self, f: &var Formatter) -> bool {
        match *self {
            Option::Some(ref v) => {
                f.write_str("Some(")
                v.fmt(f)
                f.write_str(")")
            },
            Option::None => {
                f.write_str("None")
            },
        }
        return true
    }
}
\end{lstlisting}

\subsection{Bounds on Individual Methods}

Add bounds to specific methods rather than the entire impl:

\begin{lstlisting}
// ... (illustrative - method-level where clause)
impl<T> Vec<T> {
    // Available for all Vec<T>
    pub fn len(&self) -> u32 {
        return self.length
    }

    // Only available when T: Eq
    pub fn contains(&self, value: &T) -> bool where T: Eq {
        for i in 0..self.len() {
            match self.get(i) {
                Option::Some(item) => {
                    if item.eq(value) {
                        return true
                    }
                },
                Option::None => {},
            }
        }
        return false
    }
}
\end{lstlisting}

\subsection{Const Generic Parameters}

For compile-time size parameters:

\begin{lstlisting}
struct FixedBuffer<const N: u32> {
    data: [u8; N],
    len: u32,
}

impl<const N: u32> FixedBuffer<N> {
    pub fn new() -> FixedBuffer<N> {
        return FixedBuffer {
            data: [0; N],
            len: 0,
        }
    }

    pub fn capacity(&self) -> u32 {
        return N
    }
}
\end{lstlisting}

Usage:

\begin{lstlisting}
// ... (with FixedBuffer defined above)
let small = FixedBuffer::<64>::new()   // 64-byte buffer
let large = FixedBuffer::<1024>::new() // 1KB buffer
\end{lstlisting}

\subsection{Monomorphization}

Generics in Novus are monomorphized --- the compiler generates specialized code for each concrete type used:

\begin{lstlisting}
fn max<T>(a: T, b: T) -> T where T: Ord {
    if a.cmp(&b) > 0 { return a }
    return b
}

// When you call:
let m1 = max(10, 20)
let m2 = max(5, 3)

// Compiler generates:
// fn max_i32(a: i32, b: i32) -> i32 { ... }
// fn max_u8(a: u8, b: u8) -> u8 { ... }
\end{lstlisting}

\textbf{Benefits:} Zero runtime overhead, full optimization.

\textbf{Cost:} Code size increases with each instantiation. On memory-constrained 68k systems, be mindful of instantiating generics with many different types.


\section{Common Patterns and Idioms}
\label{sec:patterns}

These patterns appear frequently in well-designed Novus code.

\subsection{The Newtype Pattern}

Wrap a type to give it distinct identity:

\begin{lstlisting}
struct Pixels {
    value: i32,
}

struct Bytes {
    value: u32,
}

struct ChipAddr {
    ptr: *u8,
}

impl Pixels {
    pub fn new(value: i32) -> Pixels {
        return Pixels { value: value }
    }

    pub fn get(&self) -> i32 {
        return self.value
    }
}
\end{lstlisting}

Now the compiler prevents mixing them up:

\begin{lstlisting}
// ... (with Pixels and Bytes defined above)
fn allocate_pixels(width: Pixels, height: Pixels) -> *u8 {
    // ...
}

fn allocate_bytes(size: Bytes) -> *u8 {
    // ...
}

let w = Pixels::new(320)
let h = Pixels::new(256)
let size = Bytes { value: 1024 }

allocate_pixels(w, h)    // OK
allocate_pixels(w, size) // Compile error! Bytes != Pixels
\end{lstlisting}

\subsubsection{Newtype for Chip RAM Addresses}

On the Amiga, chip RAM and fast RAM serve different purposes. Use newtypes to prevent bugs:

\begin{lstlisting}
struct ChipPtr {
    ptr: *u8,  // Pointer to chip RAM
}

struct FastPtr {
    ptr: *u8,  // Pointer to fast RAM
}

fn setup_copper_list(gfx: ChipPtr) {
    // Copper can only access chip RAM
    // Type system ensures we can't accidentally pass fast RAM
}

fn cache_data(buf: FastPtr) {
    // Data caching works best in fast RAM
}
\end{lstlisting}

\subsection{Extension Traits}

Add methods to types you don't own:

\begin{lstlisting}
trait I32Ext {
    fn is_even(&self) -> bool
    fn is_positive(&self) -> bool
    fn abs(&self) -> i32
}

impl I32Ext for i32 {
    fn is_even(&self) -> bool {
        return *self % 2 == 0
    }

    fn is_positive(&self) -> bool {
        return *self > 0
    }

    fn abs(&self) -> i32 {
        if *self < 0 { return -*self }
        return *self
    }
}
\end{lstlisting}

Usage:

\begin{lstlisting}
// ... (with I32Ext defined above)
let x: i32 = -42
if x.is_even() {
    println("Even!")
}
let positive = x.abs()  // 42
\end{lstlisting}

\subsection{Blanket Implementations}

Implement a trait for all types matching a bound:

\begin{lstlisting}
// ... (illustrative - every Display type gets to_string)
impl<T> ToString for T where T: Display {
    fn to_string(&self) -> Option<String> {
        match String::new() {
            Option::Some(var s) => {
                // Use Display to write to string
                self.fmt(&var s.as_formatter())
                return Option::Some(s)
            },
            Option::None => return Option::None,
        }
    }
}
\end{lstlisting}

Now any type implementing \texttt{Display} automatically gets \texttt{to\_string()}.

\subsection{Marker Traits}

Traits with no methods, used to mark capabilities:

\begin{lstlisting}
pub trait Copy {}
\end{lstlisting}

The \texttt{Copy} marker tells the compiler a type can be implicitly copied (bitwise) rather than moved. Primitives like \texttt{i32}, \texttt{u8}, and \texttt{bool} are \texttt{Copy}.

\subsection{The Builder Pattern}

For complex object construction:

\begin{lstlisting}
struct WindowConfig {
    title: *u8,
    width: u32,
    height: u32,
    borderless: bool,
}

struct WindowBuilder {
    config: WindowConfig,
}

impl WindowBuilder {
    pub fn new() -> WindowBuilder {
        return WindowBuilder {
            config: WindowConfig {
                title: "Untitled",
                width: 320,
                height: 256,
                borderless: false,
            }
        }
    }

    @chain
    pub fn title(&var self, t: *u8) {
        self.config.title = t
    }

    @chain
    pub fn size(&var self, w: u32, h: u32) {
        self.config.width = w
        self.config.height = h
    }

    @chain
    pub fn borderless(&var self) {
        self.config.borderless = true
    }

    pub fn build(&self) -> Result<WindowHandle, IntuitionError> {
        // ... create the window
    }
}
\end{lstlisting}

Usage:

\begin{lstlisting}
// ... (with WindowBuilder defined above)
var builder = WindowBuilder::new()
let window = builder
    .title("My App")
    .size(640, 480)
    .borderless()
    .build()?
\end{lstlisting}

\subsection{The Handle Pattern}

Wrap OS resources with RAII cleanup:

\begin{lstlisting}
struct FileHandle {
    lock: *FileLock,
}

impl FileHandle {
    pub fn open(path: *u8, mode: i32) -> Result<FileHandle, DosError> {
        let lock = Lock(path, mode)
        if lock == null {
            return Result::Err(dos_last_error())
        }
        return Result::Ok(FileHandle { lock: lock })
    }
}

impl Drop for FileHandle {
    fn drop(&var self) {
        if self.lock != null {
            UnLock(self.lock)
            self.lock = null
        }
    }
}
\end{lstlisting}

The file is automatically closed when the \texttt{FileHandle} goes out of scope.


\section{Summary}

This chapter covered advanced trait and generic features:

\begin{itemize}
    \item \textbf{Operator traits} --- \texttt{Add}, \texttt{Sub}, \texttt{Mul}, \texttt{Div}, \texttt{Rem}, \texttt{Neg}, \texttt{Index}, \texttt{IndexMut}
    \item \textbf{Comparison traits} --- \texttt{Eq}, \texttt{PartialOrd}, \texttt{Ord}
    \item \textbf{Conversion traits} --- \texttt{From}, \texttt{Into}, \texttt{TryFrom}, \texttt{TryInto}
    \item \textbf{Iterator traits} --- \texttt{Iterator}, \texttt{Iterable}
    \item \textbf{Numeric traits} --- \texttt{Zero}, \texttt{One}, \texttt{Bounded}
    \item \textbf{Reference traits} --- \texttt{AsRef}, \texttt{AsMut}
    \item \textbf{Advanced generics} --- multiple bounds, where clauses, const generics
    \item \textbf{Common patterns} --- newtype, extension traits, blanket implementations, builder, handle
\end{itemize}

Traits are the foundation of generic programming in Novus. Master them, and you'll write code that is both flexible and efficient --- code worthy of the Amiga.
